\documentclass[10pt,a4paper]{amsart}
\usepackage[margin=19mm]{geometry}
\usepackage{amsmath,amssymb,mathtools}
\usepackage[scr=rsfs,cal=euler]{mathalfa}
\usepackage{xfrac}
\usepackage[unicode,hidelinks,bookmarks=false]{hyperref}
\newcommand{\ie}{i.e.\ }
\newcommand{\cf}{cf.\ }
\newcommand{\resp}{respectively\ }
\newcommand{\Subsection}{\subsection}
\providecommand{\nobreakdash}{\nobreak\hbox{-}}
\setlength{\emergencystretch}{2em}
\multlinegap0pt
\usepackage{amssymb}
\usepackage{mathtools}
\usepackage{xcolor}
\usepackage[shortlabels]{enumitem}
\newtheorem{lemma}{Lemma}
\title{The Bollob\'{a}s--Nikiforov Conjecture for Complete Multipartite Graphs and Dense $K_4$-Free Graphs}
\author{Piero Giacomelli}
\date{Source: arXiv:2603.26379v3 (CC BY 4.0)}
\begin{document}
\maketitle

\noindent\emph{Source and licence.} The Bollob\'{a}s--Nikiforov Conjecture for Complete Multipartite Graphs and Dense $K_4$-Free Graphs. Version arXiv:2603.26379v3 (CC BY 4.0), distributed by arXiv under the Creative Commons Attribution 4.0 International licence, http://creativecommons.org/licenses/by/4.0/, as recorded in the arXiv licence metadata for this exact version and shown on the abstract page https://arxiv.org/abs/2603.26379v3. Full licence notice: \texttt{licenses/arxiv-2603.26379-CC-BY-4.0.txt}.\par\smallskip

\noindent\textit{Editorial scope.} Counted unit: the article's lemma lem:secular (source0.tex lines 451--464). The article's complete original proof of it is delivered with it (source0.tex lines 466--551, one \texttt{proof} environment of 86 lines). No mathematics is written, repaired or re-derived: the statement and the proof are the article's own lines, in the article's own notation, and no retained line is substituted or re-flowed.  Retained uncounted as the source-local context the proof needs: lines 408--413.  Nothing was omitted from any retained span, so the package carries no omission record.  No retained line contains a citation, so the wrapper carries no bibliography and no bibliography entry.  No retained line contains a figure environment, an image-inclusion command or drawing code, so no image is delivered and the package declares no asset dependency.  Editorial formatting only, in addition: an AMS \texttt{amsart} preamble that restates the article's own declarations and macros used by the retained lines, namely the environments \texttt{lemma} on separate counters, together with the standard packages those lines need.  The retained environments print in this document as Lemma 1 in order of appearance, whereas the article prints the same items with the numbering of its own full document.  The complete native source of the article as distributed in the arXiv e-print for this exact version is bundled unchanged as \texttt{sources/197276/source0.tex}.
\begin{lemma}[Zero eigenspace]
  \label{lem:within}
  Let $G = K_{n_1,\ldots,n_r}$ with parts $V_1,\ldots,V_r$ of sizes
  $n_1,\ldots,n_r$.  The eigenvalue $0$ of $A(G)$ has multiplicity
  exactly $n - r$.
\end{lemma}

\begin{lemma}[Secular equation and the sign of $\lambda_2$]
  \label{lem:secular}
  Let $G = K_{n_1,\ldots,n_r}$ with $r \geq 2$ and parts of sizes
  $n_1 \geq \cdots \geq n_r \geq 1$.
  The eigenvalues of $A(G)$ outside the zero eigenspace are the
  real roots of
  \begin{equation}
    \label{eq:secular}
    \sum_{i=1}^r \frac{n_i}{\lambda + n_i} \;=\; 1.
  \end{equation}
  There is exactly one positive root $\alpha_1$, and all remaining
  roots are at most $-n_r < 0$.  In particular, if $n > r$ then
  $\lambda_2(G) = 0$.
\end{lemma}

\begin{proof}
\textbf{Reduction to part-constant eigenvectors.}
Any permutation of vertices within a fixed part $V_i$ is a graph
automorphism of $K_{n_1,\ldots,n_r}$, so two vertices in the same
part have identical rows in $A(G)$.
Let $\mathbf{v}$ be an eigenvector with eigenvalue $\lambda$ that is
orthogonal to the entire zero eigenspace identified in
Lemma~\ref{lem:within}.
For each $i$, the vector $\mathbf{v}$ must be constant on $V_i$:
if it were not, the projection of $\mathbf{v}$ onto the zero eigenspace
for part $V_i$ (namely the component of $\mathbf{v}$ that is supported
on $V_i$ and sums to zero there) would be nonzero, contradicting
orthogonality.
Hence we may write $v_u = c_i$ for all $u \in V_i$.

\textbf{Deriving the secular equation.}
For $u \in V_i$, the eigenvalue equation $(A\mathbf{v})_u = \lambda c_i$
reads $\sum_{j \neq i} n_j c_j = \lambda c_i$.
Setting $s := \sum_{j=1}^r n_j c_j$, this becomes $s - n_i c_i = \lambda c_i$,
so $c_i(\lambda + n_i) = s$.
For a nontrivial eigenvector at least one $c_i$ is nonzero; if $\lambda = -n_i$
for every $i$ with $c_i \neq 0$, then $s = c_i(\lambda + n_i) = 0$ for
all $i$, forcing $s = 0$ and therefore all $c_j = s/(\lambda+n_j) = 0$
for $j$ with $\lambda \neq -n_j$; the only possibility is that $\lambda = -n_i$
for all parts $i$ with $n_i$ equal to the same value, and the eigenvectors
are differences of the constant vectors on those parts --- these are already
accounted for in the zero eigenspace when all $n_i = n_j$ (since then
$\lambda + n_j = 0$, and such a difference vector has $s = 0$).
We therefore focus on $\lambda \notin \{-n_1,\ldots,-n_r\}$ and $s \neq 0$,
giving $c_i = s/(\lambda + n_i)$.
Substituting into the definition of $s$:
\[
  s
  \;=\; \sum_{i=1}^r n_i c_i
  \;=\; s \sum_{i=1}^r \frac{n_i}{\lambda + n_i},
\]
and dividing by $s \neq 0$ yields equation~\eqref{eq:secular}.

\textbf{Root analysis.}
Define $f(\lambda) = \sum_{i=1}^r n_i/(\lambda + n_i)$.
The function $f$ is continuous and strictly decreasing on each interval
between consecutive poles, since
$f'(\lambda) = -\sum_{i=1}^r n_i/(\lambda + n_i)^2 < 0$ wherever $f$ is
defined.
The poles of $f$ occur at the values $\{-n_i : 1 \leq i \leq r\}$;
let the distinct values in this set be
$-p_1 < -p_2 < \cdots < -p_s \leq -1 < 0$,
where $p_1 > p_2 > \cdots > p_s \geq 1$ are the distinct part sizes
and $s \leq r$.

\textit{One positive root.}
On $(0,+\infty)$, every term $n_i/(\lambda + n_i)$ is positive and
decreasing to $0$, so $f$ decreases strictly from $f(0) = r \geq 2 > 1$
to $\lim_{\lambda \to +\infty} f(\lambda) = 0 < 1$.
The intermediate value theorem gives a unique root $\alpha_1 \in (0,+\infty)$.

\textit{No root in $(-p_s, 0)$.}
On the interval $(-p_s, 0)$, every denominator $\lambda + n_i$ satisfies
$\lambda + n_i \geq \lambda + p_s > 0$, so $f(\lambda) \geq 0$.
Moreover, as $\lambda \to (-p_s)^+$, the term $p_s/(\lambda + p_s)$
(summed over all parts with $n_i = p_s$) diverges to $+\infty$, so
$\lim_{\lambda \to (-p_s)^+} f(\lambda) = +\infty$.
Since $f$ is strictly decreasing on $(-p_s,0)$ and $f(0) = r > 1$,
we conclude $f(\lambda) > r > 1$ for all $\lambda \in (-p_s, 0)$,
hence no root lies in this interval.

\textit{All remaining roots are at most $-p_s < 0$.}
On each inter-pole interval $(-p_{k+1},-p_k)$ for $k = 1,\ldots,s-1$:
as $\lambda \to (-p_{k+1})^+$, the terms with $n_i = p_{k+1}$ give
$f \to +\infty$, and as $\lambda \to (-p_k)^-$, the terms with
$n_i = p_k$ give $f \to -\infty$.
By the intermediate value theorem and strict monotonicity, there is
exactly one root in each such interval.
No root lies in $(-\infty, -p_1)$ because $f(\lambda) < 0 < 1$ there
(all denominators are negative and all numerators positive).
Including the poles themselves: if $\lambda = -n_j$ for some $j$ with
$\lambda + n_i \neq 0$ for all $i \neq j$, then $f(\lambda)$ is undefined,
so $\lambda$ is not a root of \eqref{eq:secular}.

\textbf{Conclusion.}
All roots of \eqref{eq:secular} outside the zero eigenspace are:
the unique positive root $\alpha_1$, and roots at most $-p_s \leq -1 < 0$.
Together with the zero eigenspace of dimension $n - r \geq 1$
(when $n > r$), the second largest eigenvalue of $A(G)$ is
$\lambda_2(G) = 0$.
\end{proof}

\end{document}
